% Options for packages loaded elsewhere
\PassOptionsToPackage{unicode}{hyperref}
\PassOptionsToPackage{hyphens}{url}
\documentclass[
  11pt,
]{article}
\usepackage{xcolor}
\usepackage[margin=1in]{geometry}
\usepackage{amsmath,amssymb}
\setcounter{secnumdepth}{-\maxdimen} % remove section numbering
\usepackage{iftex}
\ifPDFTeX
  \usepackage[T1]{fontenc}
  \usepackage[utf8]{inputenc}
  \usepackage{textcomp} % provide euro and other symbols
\else % if luatex or xetex
  \usepackage{unicode-math} % this also loads fontspec
  \defaultfontfeatures{Scale=MatchLowercase}
  \defaultfontfeatures[\rmfamily]{Ligatures=TeX,Scale=1}
\fi
\usepackage{lmodern}
\ifPDFTeX\else
  % xetex/luatex font selection
\fi
% Use upquote if available, for straight quotes in verbatim environments
\IfFileExists{upquote.sty}{\usepackage{upquote}}{}
\IfFileExists{microtype.sty}{% use microtype if available
  \usepackage[]{microtype}
  \UseMicrotypeSet[protrusion]{basicmath} % disable protrusion for tt fonts
}{}
\makeatletter
\@ifundefined{KOMAClassName}{% if non-KOMA class
  \IfFileExists{parskip.sty}{%
    \usepackage{parskip}
  }{% else
    \setlength{\parindent}{0pt}
    \setlength{\parskip}{6pt plus 2pt minus 1pt}}
}{% if KOMA class
  \KOMAoptions{parskip=half}}
\makeatother
\usepackage{longtable,booktabs,array}
\newcounter{none} % for unnumbered tables
\usepackage{calc} % for calculating minipage widths
% Correct order of tables after \paragraph or \subparagraph
\usepackage{etoolbox}
\makeatletter
\patchcmd\longtable{\par}{\if@noskipsec\mbox{}\fi\par}{}{}
\makeatother
% Allow footnotes in longtable head/foot
\IfFileExists{footnotehyper.sty}{\usepackage{footnotehyper}}{\usepackage{footnote}}
\makesavenoteenv{longtable}
\setlength{\emergencystretch}{3em} % prevent overfull lines
\providecommand{\tightlist}{%
  \setlength{\itemsep}{0pt}\setlength{\parskip}{0pt}}
\usepackage{bookmark}
\IfFileExists{xurl.sty}{\usepackage{xurl}}{} % add URL line breaks if available
\urlstyle{same}
\hypersetup{
  hidelinks,
  pdfcreator={LaTeX via pandoc}}

\author{}
\date{}


\usepackage{graphicx}
\usepackage{float}
\begin{document}

\section{Teaching a Coding Agent to Design DNA Origami: What Worked,
What Didn't, and What It Took to Get
There}\label{teaching-a-coding-agent-to-design-dna-origami-what-worked-what-didnt-and-what-it-took-to-get-there}

Daniel Fu, {[}Co-authors TBD{]}

\emph{Target: ISNSCE DNA 32 --- Track B (Extended Abstract) or Poster}

\subsection{Abstract}\label{abstract}

We report on our experience using a general-purpose coding agent (Claude
Code) to design DNA origami structures in cadnano. We describe two
approaches --- an embedded tool-using LLM that failed, and a coding
agent with source code access that succeeded --- and trace the path from
initial failure to a working parametric pipeline for a 2-layer
rectangular origami with a cavity. The key finding is that the agent's
bottleneck was not code generation but domain knowledge: each class of
design error required explicit correction from the human designer before
the agent could proceed. We detail the specific errors, the domain
knowledge that unblocked each one, and the capabilities that emerged
after unblocking.

\subsection{1. The Problem}\label{the-problem}

Modifying a DNA origami design is slow. When a cavity is too narrow by 5
nm, the designer must re-route the scaffold, re-place staples, re-apply
edge treatments, and re-verify the structure. This process takes hours
and must be repeated for every design iteration. We asked: can an AI
agent handle this redesign automatically?

\subsection{2. First Attempt: Embed an LLM Inside cadnano
(Failed)}\label{first-attempt-embed-an-llm-inside-cadnano-failed}

We built an agent dialog into the cadnano GUI. The user types a natural
language command; an LLM parses the intent and calls Python methods that
modify the design. We implemented 85 methods covering helix creation,
strand manipulation, crossover placement, and staple breaking.

This failed on anything beyond single-step operations. Scaffold routing
--- the task of placing crossovers to create one continuous scaffold
loop --- requires reasoning about helix parity, crossover position
tables, and routing topology simultaneously. Even GPT-class models
achieved 0\% success. We attempted reinforcement learning with a design
verifier as reward signal, but the local model (Qwen 1.7B) could not
discover correct action sequences through exploration.

\textbf{What we learned:} The 85-method API encodes necessary but
insufficient domain knowledge. The model must compose methods into
correct sequences without understanding why each step matters. The API
is a keyhole --- the model can see the actions but not the mechanism.

\subsection{3. Second Attempt: A Coding Agent with Source Code Access
(Worked, Then Didn't, Then
Did)}\label{second-attempt-a-coding-agent-with-source-code-access-worked-then-didnt-then-did}

We switched to a coding agent (Claude Code) that reads cadnano's source
code and writes Python scripts. This immediately solved several problems
the embedded agent could not:

\begin{itemize}
\tightlist
\item
  When \texttt{createXover} broke scaffold continuity by splitting
  strands, the agent found \texttt{setConnection3p/5p} by reading the
  strand model internals.
\item
  When the autobreak plugin failed to import, the agent traced the
  module system and used \texttt{importlib} to bypass a legacy
  \texttt{\_\_init\_\_.py}.
\item
  The agent integrated cadnano with tacoxDNA (format conversion) and
  oxDNA (molecular dynamics), creating an end-to-end pipeline no single
  tool provides.
\end{itemize}

\textbf{First success:} A rectangular origami with the m13mp18 scaffold
(7,249 nt). 24 helices, serpentine routing, autoStaple, autoBreak, twist
correction, tacoxDNA conversion, oxDNA relaxation. This is trivial for a
human designer --- but it proved the pipeline worked.

\subsection{4. Scaling Up: Where Domain Knowledge Became the
Bottleneck}\label{scaling-up-where-domain-knowledge-became-the-bottleneck}

We then attempted a 2-layer rectangle with a rectangular cavity using
the p8064 scaffold (8,064 nt). This is where the agent's limitations
became clear. Each step forward required the human designer to diagnose
and correct a specific class of error.

\subsubsection{Error 1: Wrong grid
layout}\label{error-1-wrong-grid-layout}

The agent built a 4$\times$7 grid (4 rows, 7 columns). A 2-layer rectangle
should be 2$\times$12 (2 grid rows, 12 columns). The agent did not know that
``2-layer'' means 2 grid rows on the honeycomb lattice, because this
mapping depends on understanding how honeycomb parity creates the
physical layer structure. \textbf{Unblocked by:} the PI explaining the
grid-to-layer correspondence.

\subsubsection{Error 2: Serpentine routing instead of dense
crossovers}\label{error-2-serpentine-routing-instead-of-dense-crossovers}

The agent routed the scaffold with one crossover per helix pair
(serpentine). Real DNA origami uses multiple double crossovers per pair
--- the scaffold weaves back and forth between adjacent helices. This is
a fundamental design convention the agent had no way to discover from
the source code alone. \textbf{Unblocked by:} the PI providing a
hand-designed template (\texttt{2x12\_rectangle\_cavity.json}) showing
the correct crossover pattern.

\subsubsection{Error 3: Cavity in the wrong
dimension}\label{error-3-cavity-in-the-wrong-dimension}

The agent placed the cavity by removing helices from the XY
cross-section (changing the grid). The cavity should be in the Z
direction --- a gap along the helix length axis, visible in the side
view. \textbf{Unblocked by:} the PI specifying that the cavity cuts
through the helix length, not the helix arrangement.

\subsubsection{Error 4: Scaffold/staple
confusion}\label{error-4-scaffoldstaple-confusion}

When building from scratch, the agent placed scaffold strands on the
staple strand set and vice versa. \textbf{Unblocked by:} the PI
inspecting the design in cadnano and identifying the swap.

\subsubsection{Error 5: Scaling broke file
format}\label{error-5-scaling-broke-file-format}

When scaling the PI's template from 252 to 420 bp, the agent's scaling
script left 5,036 stray staple entries with broken crossover references.
cadnano crashed on load. \textbf{Unblocked by:} tracing the crash to
\texttt{legacydecoder.py} line 205, identifying the staple crossover
references as the cause, and clearing them.

\subsubsection{Error 6: Cavity gaps filled
incorrectly}\label{error-6-cavity-gaps-filled-incorrectly}

When extending the design, the agent filled the cavity gap with
continuous scaffold --- destroying the cavity. \textbf{Unblocked by:}
the PI manually clearing positions 75--337 on the cavity helices and
returning the corrected file.

\subsubsection{Error 7: Crossover
misalignment}\label{error-7-crossover-misalignment}

The agent's scaling logic moved crossover positions independently per
helix, creating skewed crossovers (e.g., H16{[}337{]} vs H17{[}338{]}).
Both sides of a crossover must be at the same base pair index.
\textbf{Unblocked by:} the PI aligning the crossovers at index 338 on
both helices.

\subsubsection{Error 8: Cavity edges not aligned between
layers}\label{error-8-cavity-edges-not-aligned-between-layers}

The cavity right edge was at position 270 in one layer and 338 in the
other. The PI requested they be aligned. The nearest valid honeycomb
position on each layer differs by 2 bp (340 vs 338) due to different
crossover direction offsets. \textbf{Resolved by:} the agent computing
valid crossover positions for each honeycomb direction and selecting the
closest match.

\subsection{5. What the Agent Could Do After Each
Unblocking}\label{what-the-agent-could-do-after-each-unblocking}

After each correction, the agent incorporated the lesson and did not
repeat that specific error class. The progression:

\begin{enumerate}
\def\labelenumi{\arabic{enumi}.}
\tightlist
\item
  \textbf{After Error 1} (grid layout): Correctly built 2$\times$12 grids from
  then on.
\item
  \textbf{After Error 2} (dense crossovers): Used the PI's template as
  the basis for all subsequent designs instead of building from scratch.
\item
  \textbf{After Error 3} (cavity orientation): All cavity designs used
  Z-direction gaps.
\item
  \textbf{After Errors 5--7} (scaling): Developed a 4-step scaling
  process that the PI validated step-by-step:

  \begin{itemize}
  \tightlist
  \item
    Step 1: Remove staples (keep scaffold only). Verified loads.
  \item
    Step 2: Extend helix arrays and shift right-side crossovers.
    Verified loads.
  \item
    Step 3: Move midseam crossovers to center. Verified loads.
  \item
    Step 4: Expand cavity by moving boundary crossovers. Verified loads.
  \end{itemize}
\end{enumerate}

The final result: a 2$\times$12 rectangular origami with an aligned cavity, 1
continuous scaffold loop (7,688 bp), scaled from the PI's 252 bp
template to 420 bp helices. Each step produced a valid cadnano JSON that
loads without errors.

\subsection{6. Beyond the Template: Autonomous Design at
Scale}\label{beyond-the-template-autonomous-design-at-scale}

After the PI provided sufficient guidance through 10 failure modes, the
agent achieved autonomous design capability. Without further human
intervention, the agent:

\begin{enumerate}
\def\labelenumi{\arabic{enumi}.}
\tightlist
\item
  \textbf{Extended the template to arbitrary widths} (2$\times$14, 2$\times$16, 2$\times$18,
  2$\times$20) by adding column pairs to both sides of the 2$\times$12 template,
  maintaining 1 scaffold oligo at each step.
\item
  \textbf{Designed to specification:} Given the constraint ``20 nm $\times$ 40
  nm cavity, centered, 6-helix padding, scaffold $\leq$ 8,064 bp,'' the agent
  independently selected a 2$\times$20 grid (40 helices $\times$ 252 bp), computed the
  cavity dimensions (8 columns $\times$ 117 bp gap), and built the complete
  design --- 7,828 bp scaffold, 1 oligo.
\item
  \textbf{Executed targeted edits from natural language prompts:} When
  told ``the crossovers at H5-H6 and H33-H34 are too close to the cavity
  edge --- move them away,'' the agent identified the positions
  ({[}64,65{]} and {[}183,184{]}, both 4 bp from the cavity), found
  valid alternative lattice positions ({[}43,44{]} and {[}204,205{]}, 25
  bp clearance), and moved them while preserving the single scaffold
  oligo (Figure 5).
\item
  \textbf{Ran the full pipeline autonomously:} autoStaple (222 staples)
  $\to$ autoBreak (all3) $\to$ tacoxDNA (15,656 nt) $\to$ oxDNA PACE GPU relaxation.
\end{enumerate}

\subsection{7. Agents as a Substrate for Shared Design
Intelligence}\label{agents-as-a-substrate-for-shared-design-intelligence}

The key insight of this work is not that AI can design DNA origami ---
it cannot, at least not independently. The insight is that the
\emph{failures} are as valuable as the successes. Each of the 10 failure
modes we documented represents a piece of domain knowledge that:

\begin{enumerate}
\def\labelenumi{\arabic{enumi}.}
\tightlist
\item
  \textbf{Is invisible in the source code.} The rule ``LOW position =
  crossover IN, HIGH = crossover OUT'' cannot be deduced from cadnano's
  codebase. It must be taught by a human who has internalized years of
  design experience.
\item
  \textbf{Is immediately reusable.} Once the agent learns a rule, it
  never makes that class of error again. The learning is encoded in
  verifier functions, lessons files, and pipeline logic that persist
  across sessions.
\item
  \textbf{Is distributable.} The verifier
  (\texttt{cadnano\_verifier.py}) catches all 10 failure modes
  automatically. A new user --- or a new AI agent --- can skip our
  failures entirely. The failure catalog (\texttt{failure\_analysis.md})
  explains each one with symptom, root cause, fix, and verifier check.
\end{enumerate}

This creates a new paradigm: \textbf{the agent is not the designer, but
the substrate on which shared design intelligence accumulates.} Each
PI-agent interaction produces not just a design, but reusable tools that
make the next design faster. The 10 failures we encountered over two
weeks become the 10 checks that a future user's agent runs in seconds.

\subsubsection{What Can Be Shared and
Replicated}\label{what-can-be-shared-and-replicated}

The following artifacts are distributable to any DNA origami researcher
working with cadnano:

{\def\LTcaptype{none} % do not increment counter
\begin{longtable}[]{@{}
  >{\raggedright\arraybackslash}p{(\linewidth - 4\tabcolsep) * \real{0.3125}}
  >{\raggedright\arraybackslash}p{(\linewidth - 4\tabcolsep) * \real{0.2812}}
  >{\raggedright\arraybackslash}p{(\linewidth - 4\tabcolsep) * \real{0.4062}}@{}}
\toprule\noalign{}
\begin{minipage}[b]{\linewidth}\raggedright
Artifact
\end{minipage} & \begin{minipage}[b]{\linewidth}\raggedright
Purpose
\end{minipage} & \begin{minipage}[b]{\linewidth}\raggedright
Reusability
\end{minipage} \\
\midrule\noalign{}
\endhead
\bottomrule\noalign{}
\endlastfoot
\texttt{cadnano\_verifier.py} & Pre-flight design validation & Run on
any cadnano JSON before committing to synthesis \\
\texttt{failure\_analysis.md} & Documented failure modes & Read before
starting AI-assisted design; avoid 10 known pitfalls \\
\texttt{cavity\_variant\_sweep.py} & Parametric cavity pipeline & Change
gap size/width $\to$ design recomputes automatically \\
Template extension functions & Add columns to existing designs & Extend
any 2$\times$N template to 2$\times$(N+2) with 1 scaffold oligo \\
Crossover move pattern & LOW=IN, HIGH=OUT template & Apply to any
targeted crossover edit on any honeycomb design \\
Renumbering procedure & JSON helix renumbering & Essential for any
template extension; ensures parity correctness \\
\end{longtable}
}

The pipeline from PI's template to verified design to oxDNA simulation
is end-to-end: no manual cadnano GUI interaction required. A researcher
can describe a cavity design in terms of physical dimensions (nm $\times$ nm,
scaffold length, padding) and receive a validated JSON, stapled design,
and oxDNA files ready for PACE submission.

\subsubsection{The Threshold We Crossed}\label{the-threshold-we-crossed}

Early in this project, the agent produced ``cool demos'' --- a flat
sheet, a simple rectangle. These demonstrated capability but had no
transferable value. The threshold from demo to knowledge transfer was
crossed when the agent began \emph{building its own verification tools
from its own failures.} The verifier was not designed top-down; it
emerged bottom-up from 10 specific debugging sessions. This is the
mechanism by which an agent becomes a substrate for shared intelligence:
each failure produces a check that prevents the same failure for every
future user.

\subsection{8. The Final Demonstration: One-Shot Targeted
Edits}\label{the-final-demonstration-one-shot-targeted-edits}

The strongest evidence of transferred design intelligence came in the
final session. The PI issued three successive modification requests on a
2$\times$22 integrin cavity design, and the agent executed each one correctly
on the first attempt:

\subsubsection{Edit 1: ``Shrink the
structure''}\label{edit-1-shrink-the-structure}

The 2$\times$22 design was wider than needed. Rather than truncating helix
arrays (which would destroy edge crossovers), the agent recognized the
correct approach: \textbf{move the edge crossovers inward} by one step
(from positions 246/242 to 225/221), then clear scaffold data beyond the
new edges. This preserved all crossover connectivity while reducing the
effective helix length from 252 bp to 225 bp. Result: 7,412 bp scaffold,
1 oligo.

\subsubsection{Edit 2: ``Move the full crossovers away from the cavity
edge, and center the
cavity''}\label{edit-2-move-the-full-crossovers-away-from-the-cavity-edge-and-center-the-cavity}

The agent identified 5 inter-pair crossovers that were only 3--5 bp from
cavity boundary half-crossovers (H5-H6 at {[}64,65{]}, H13-H14 at
{[}53,54{]}, H15-H16 at {[}64,65{]}, H27-H28 at {[}183,184{]}, H37-H38
at {[}183,184{]}). It moved each to a safe lattice position 26+ bp from
the cavity. Then it recomputed centered cavity boundaries: R12 gap
{[}59--171{]} with 54 bp on each side, R13 gap {[}55--168{]} with 53 bp
on each side. Result: 7,432 bp scaffold, 1 oligo, cavity centered to
within 1 bp.

\subsubsection{Edit 3: ``Now autoStaple and autoBreak
it''}\label{edit-3-now-autostaple-and-autobreak-it}

The agent ran autoStaple (82 initial staples) and autoBreak with
minStapleLegLen=3 (220+ broken staples, 18--279 bp range, preserving
full crossovers). Result: complete stapled design ready for tacoxDNA
conversion.

\textbf{Why this matters:} Each edit required combining multiple pieces
of domain knowledge that were transferred across previous sessions ---
crossover lattice positions, parity-dependent scaffold direction, cavity
boundary half-crossover placement, honeycomb valid offsets. None of
these were re-taught. The agent drew on lessons encoded in
\texttt{tasks/lessons.md}, verifier checks in
\texttt{cadnano\_verifier.py}, and pipeline functions in
\texttt{cavity\_variant\_sweep.py}. This is the mechanism of
transferable design intelligence: the PI's corrections accumulate as
reusable code, and the agent applies them automatically in new contexts.


\section{Figures}

\begin{figure}[H]
\centering
\includegraphics[width=0.9\textwidth]{../figures/fig1_architecture.png}
\caption{Architecture comparison: embedded tool-using LLM (failed) vs.\ coding agent with source code access (succeeded).}
\label{fig:figure1}
\end{figure}

\begin{figure}[H]
\centering
\includegraphics[width=0.9\textwidth]{../figures/fig2_pi_template.png}
\caption{The PI's template (\texttt{2x12\_rectangle\_cavity.json}) showing correct scaffold routing with dense crossovers and cavity gap.}
\label{fig:figure2}
\end{figure}

\begin{figure}[H]
\centering
\includegraphics[width=0.9\textwidth]{../figures/fig3_cavity_design.png}
\caption{The 4-step scaling process: (a) scaffold only, (b) extended right side, (c) midseam moved, (d) cavity expanded and aligned between layers.}
\label{fig:figure3}
\end{figure}

\begin{figure}[H]
\centering
\includegraphics[width=0.9\textwidth]{../figures/fig4_parametric_sweep.png}
\caption{Scaffold routing visualization of the final scaled design (420 bp, 7,688 bp scaffold, 1 oligo).}
\label{fig:figure4}
\end{figure}

\begin{figure}[H]
\centering
  \begin{minipage}{0.48\textwidth}\centering
    \includegraphics[width=\linewidth]{../figures/fig5a_targeted_before.png}
  \end{minipage}\hfill
  \begin{minipage}{0.48\textwidth}\centering
    \includegraphics[width=\linewidth]{../figures/fig5b_targeted_after.png}
  \end{minipage}\hfill
\caption{Targeted crossover edit. Left: before (scaffold-only). Right: after. H5--H6 moved from [64,65] to [43,44]; H33--H34 moved from [183,184] to [204,205].}
\label{fig:figure5}
\end{figure}

\begin{figure}[H]
\centering
\includegraphics[width=0.9\textwidth]{../figures/fig6_oligo_journey.png}
\caption{Scaffold oligo count journey: 65 $\to$ 54 $\to$ 10 $\to$ 8 $\to$ 4 $\to$ 1.}
\label{fig:figure6}
\end{figure}

\begin{figure}[H]
\centering
\includegraphics[width=0.9\textwidth]{../figures/fig7_verifier_comparison.png}
\caption{\texttt{cadnano\_verifier.py} output: before (3 failures, 8 oligos) vs.\ after (all pass, 1 oligo).}
\label{fig:figure7}
\end{figure}

\begin{figure}[H]
\centering
\includegraphics[width=0.9\textwidth]{../figures/fig8_cavities_oxdna.png}
\caption{All three parametric cavity variants (20/30/40 nm gap) after oxDNA production relaxation (20M MD steps).}
\label{fig:figure8}
\end{figure}

\begin{figure}[H]
\centering
\includegraphics[width=0.9\textwidth]{../figures/fig9_targeted_edits.png}
\caption{One-shot targeted edits from natural language prompts on the 2$\times$22 integrin cavity design.}
\label{fig:figure9}
\end{figure}

\begin{figure}[H]
\centering
\includegraphics[width=0.9\textwidth]{../figures/fig_failure_table.png}
\caption{Summary of failure modes, domain knowledge required, and resolution.}
\label{fig:failuretable}
\end{figure}

\end{document}
